\section{The exclusion of sporadic normal cores}
\label{sec:sporadic-core}

\paragraph{Status and input.}
The proofs and exact finite certificates in this section are new.
The complete mathematical input consists of the ordinary character
tables of the twenty-six universal sporadic covers and twelve outer
extensions, with their stored power maps and subgroup fusions, in
the author-maintained Character Table Library, version 1.3.11.
Every one of these thirty-eight tables identifies its origin as the
published \emph{ATLAS of Finite Groups}; the exact table identifiers,
source strings, and complete data are retained in
\texttt{character\_inputs.json} and \texttt{coverage.json}.
We use these finite tables and their identification with the named
groups as published mathematical input, not as groups reconstructed
from orthogonality alone.  The arithmetic certificates below have
been executed.  Both the parent and a sibling agent independently
read the complete proof and programs and returned PASS.  The parent
reran the standalone verifier; the sibling additionally rebuilt the
full calculation using independent cyclotomic arithmetic, restriction
inner products, and Ramanujan-sum Fourier multiplicities.  Both
reconstructions passed all thirty-eight tables, 3035 selected rows,
and 473972 nonscalar entries.  Exact review scope and source-input
limitations are retained in the check register and peer-audit report.

For a finite-order operator on an $e$-dimensional complex space put
\[
 \nu(g)=e-\max_\lambda\dim\ker(g-\lambda).
\]
This invariant is unchanged by multiplication by a nonzero scalar.
A nonidentity scalar has $\nu(g)=0$ but actual moving rank $e$.

\subsection*{D-SPOR-01. Complete finite character input}

\begin{lemma}\label{D-SPOR-01}
Let $S$ be any sporadic simple group and $U$ its universal perfect
central cover.  Every nonlinear irreducible character of $U$ has
degree at least six.  In every such representation of degree $e$,
every nonscalar element has $\nu\ge e/4$.

If $\operatorname{Out}(S)$ is nontrivial, it has order two.  Let
$A$ be the cover-by-outer-involution extension in Table~\ref{tab:sporadic-coverage}.
For every irreducible character of $A$ whose restriction to $U$ is
nonlinear and irreducible, the same bound holds for every nonscalar
element of $A$, including elements outside $U$.
\end{lemma}

\begin{table}[ht]
\centering\small
\begin{tabular}{c|c|c|c@{\qquad}c|c|c|c}
$S$&$|Z(U)|$&$|\operatorname{Out}(S)|$&$A$&
$S$&$|Z(U)|$&$|\operatorname{Out}(S)|$&$A$\\ \hline
$B$&2&1&---&$J_4$&1&1&---\\
$Co_1$&2&1&---&$Ly$&1&1&---\\
$Co_2$&1&1&---&$M$&1&1&---\\
$Co_3$&1&1&---&$M_{11}$&1&1&---\\
$Fi_{24}'$&3&2&$3.F3+.2$&$M_{12}$&2&2&$2.M12.2$\\
$Fi_{22}$&6&2&$6.Fi22.2$&$M_{22}$&12&2&$12.M22.2$\\
$Fi_{23}$&1&1&---&$M_{23}$&1&1&---\\
$HN$&1&2&$HN.2$&$M_{24}$&1&1&---\\
$HS$&2&2&$2.HS.2$&$McL$&3&2&$3.McL.2$\\
$He$&1&2&$He.2$&$O'N$&3&2&$3.ON.2$\\
$J_1$&1&1&---&$Ru$&2&1&---\\
$J_2$&2&2&$2.J2.2$&$Suz$&6&2&$6.Suz.2$\\
$J_3$&3&2&$3.J3.2$&$Th$&1&1&---
\end{tabular}
\caption{Complete sporadic coverage.  The rightmost entries in each
half are exact CTblLib identifiers.  The cover identifier is obtained
by prefixing the displayed central order when it exceeds one.
CTblLib writes $Fi_{24}'$ as $F3+$.}
\label{tab:sporadic-coverage}
\end{table}

\begin{proof}
The group and multiplier identifications in the table are the
\emph{ATLAS} data \cite{SPOR:Atlas}, accessed through
\texttt{ExtensionInfoCharacterTable} and the complete
\texttt{IsSporadicSimple} list of \cite{SPOR:CTblLib}.
The latter attribute records the Schur multiplier and outer
automorphism group, and specifies the corresponding cover and
automorphism-group identifiers; its precise manual locator is
\cite[\S3.7, entry \texttt{ExtensionInfoCharacterTable}]{SPOR:CTblLib}.
All twelve indicated extension tables and their actual index-two
subgroup fusions from $U$ are present.  No group is omitted because
its order or representation degree is large: the Baby Monster and
Monster tables are included.

Here is the finite proof of the asserted spectral bounds, including
the exact test performed on every entry.  A character value $z$ is
encoded as an integer, or as an integral polynomial in
$\zeta_n=\exp(2\pi i/n)$.  Compute $a=z\overline z$ exactly.
If $a=e^2$, the operator is scalar: equality in the triangle
inequality for a sum of $e$ unit complex eigenvalues forces them
all to coincide.  Conversely a scalar has $a=e^2$.  Such classes
are discarded only for the assertion about $\nu$.

For a nonscalar entry write $a=\sum b_j\zeta_n^j$.  The rational
integer $B=\sum|b_j|$ is an upper bound for $|z|^2$.  Whenever
$4B\le e^2$, the trace of the operator after rescaling its largest
eigenspace to eigenvalue one gives
\[
 e-2\nu\le |z|\le e/2,
 \qquad\text{and therefore}\qquad \nu\ge e/4.       \tag{SPOR.1}
\]
When that sufficient triangle bound fails, use the character values
on powers and the exact Fourier formula.  For a class represented
by an element of order $n$, the multiplicity of $\zeta_n^k$ is
\[
 m_k=\frac1n\sum_{r=0}^{n-1}\chi(g^r)\zeta_n^{-kr}.
                                                               \tag{SPOR.2}
\]
Every number in (SPOR.2) is computed in a cyclotomic field with
integer coefficients.  The certificate verifies that each $m_k$
is a nonnegative integer, that $\sum m_k=e$, and that
$4(e-\max m_k)\ge e$.  This formula is just the orthogonality
formula for the characters of the cyclic subgroup; no numerical
approximation or eigenvalue recognition is involved.

For an extension table, retain exactly the rows whose restriction
along the stored subgroup fusion has inner product one with itself
on $U$, and whose degree exceeds one.  These are precisely the
characters specified in the lemma.  The executed GAP certificate
\texttt{scan\_quarter\_bound.g} checks all twenty-six cover tables
and twelve extension tables.  It covers $3035$ selected rows and
$473972$ nonscalar entries.  Its triangle test needs (SPOR.2) only
for $24$ entries in the table $2.J2$, and all pass.  They occur in
rows $2,3$ of degree $14$ and rows $22,23$ of degree $6$; their
computed codimensions are respectively $10$ and $4$.  The minimum
nonlinear cover degree is six.  These are complete scans of the
rows and columns, not lists of favorable conjugacy classes.

The separate program \texttt{export\_inputs.g} exports all the
complete tables, not merely the selected rows or exceptional values.
It also verifies exact row orthogonality, the degree-squared sum,
and the index-two subgroup data.  The independent arithmetic
implementation \texttt{verify\_quarter\_bound.py} calls no GAP
code.  It computes in $\mathbf Z[x]/(\Phi_n(x))$, independently
selects the extension rows by equality of the restricted character
with a row of the complete cover table, and reconstructs the
Fourier multiplicities from prime power maps.  Its different
polynomial basis needs only twelve Fourier fallbacks; it verifies
the same $3035$ rows and $473972$ entries and again finds minimum
degree six and no failed inequality.  The exact input SHA-256 is
\texttt{c2268f05043ed29a9048d5fe7d55b1a096e5c8615fa049b3438cca088f703edb}.

Thus the elementary test (SPOR.1)--(SPOR.2), applied exhaustively
to the stated published input, proves the finite lemma.  The
library notes a table-automorphism ambiguity for a power map in
$3.McL$; no Fourier fallback uses that table, so this ambiguity
does not affect the proof.  The exceptional Fourier checks all
use the stored $2.J2$ power maps.
\end{proof}

\subsection*{D-SPOR-02. All scalar and outer normalizers}

\begin{proposition}\label{D-SPOR-02}
Let $N\subset\operatorname{GL}(W)$ be a finite quasisimple group
acting faithfully and irreducibly, with sporadic simple quotient,
and put $a=\dim W$.  Every nonscalar element of its linear
normalizer has largest-eigenspace codimension at least $a/4$.
The normalizing element may have any scalar phase.  No equality
of character fields or model over an imaginary quadratic field
is required.
\end{proposition}

\begin{proof}
Let $S=N/Z(N)$ and let $U\to N$ be the quotient of the universal
perfect central cover.  The given representation becomes a
nonlinear irreducible representation $\rho$ of $U$.
The centralizer of $N$ in $\operatorname{GL}(W)$ is scalar, by
Schur's lemma.  An automorphism of the perfect group $N$ that
induces the identity on $S$ is the identity: its discrepancy from
the identity is a homomorphism from $N$ to the central abelian
group $Z(N)$.  Thus the linear normalizer modulo scalars embeds
in $\operatorname{Aut}(S)$, and its quotient by the projective
image of $N$ embeds in $\operatorname{Out}(S)$.

If this outer image is trivial, every normalizer element is a
scalar multiple of an element of $N$, and D-SPOR-01 applies.
Otherwise $\operatorname{Out}(S)=C_2$.  Fix the actual finite
extension $A=U.2$ used in D-SPOR-01 and choose $t\in A\setminus U$.
Write $\alpha$ for conjugation by $t$ on $U$ and $u_0=t^2\in U$.
This argument does not require $t^2=1$ or a split extension.

Let $g$ be a normalizing matrix inducing an outer automorphism.
After multiplying $g$ by an element of $N$, its induced
automorphism of $S$ agrees with that of $\alpha$.  Automorphisms
of $S$ lift uniquely to its universal perfect cover: this follows
from the universal property by pulling back the central extension,
and uniqueness also follows because a discrepancy is a homomorphism
from the perfect group $U$ to its center.  Consequently $\alpha$
preserves the kernel of $\rho$, and
\[
             g\rho(u)g^{-1}=\rho(\alpha(u))\qquad(u\in U)
                                                               \tag{SPOR.3}
\]
after the indicated inner adjustment.  In particular the character
of $\rho$ is $\alpha$-invariant.

By (SPOR.3), $g^2\rho(u_0)^{-1}$ centralizes $\rho(U)$ and is
a nonzero scalar.  Multiply $g$ by a scalar square root of its
inverse to obtain $T$ with $T^2=\rho(u_0)$.  Equation (SPOR.3)
and this square relation give a representation of the actual
index-two extension $A$, extending $\rho$ and sending $t$ to $T$.
It is irreducible because its restriction to $U$ is irreducible.
Therefore its character is one of the extension rows selected in
D-SPOR-01.  The original normalizing element is a scalar multiple
of an element of this represented $A$, up to the inner adjustment
already made.  The bound of D-SPOR-01 is unchanged by such scalar
multiplication.  This proves the proposition.

The argument also explains why no isoclinic version of $U.2$ is
missing: an alternative square relation changes only the scalar
normalization of its intertwiner.  We use the actual fixed table
$A$ and extend to it directly.  Characters that are exchanged by
the outer automorphism cannot occur as irreducible restrictions
of a normalizing representation and are correctly omitted.
\end{proof}

\subsection*{D-SPOR-03. Tensor permutations and the sixfold exclusion}

\begin{theorem}\label{D-SPOR-03}
Suppose a complex inducing representation for a finite-cover
sixfold source has faithful finite image $J\subset\operatorname{GL}(V)$.
Suppose $J$ has a normal subgroup $E$ that acts irreducibly on $V$
and is a central product of quasisimple groups with sporadic simple
quotients.  Then the actual six-dimensional multiplicity variation
cannot have period-map rank nine.

This excludes sporadic cores in D-BI-19, including arbitrary outer
automorphisms and permutations of quasisimple tensor factors.  It
does not require a nonsplit polarization and permits an inducing
character field larger than the original quadratic field.
\end{theorem}

\begin{proof}
Write $E=E_1\cdots E_t$ with commuting quasisimple factors.
Irreducibility gives
\[
                   V=V_1\otimes\cdots\otimes V_t,
                   \qquad a_i=\dim V_i\ge6.                 \tag{SPOR.4}
\]
Each factor representation is irreducible and faithful: an element
acting as the identity on $V_i$ acts as the identity on $V$, and
$E_i$ is an actual matrix subgroup.  The factor degrees are at
least six by D-SPOR-01.

An automorphism of $E$ permutes its quasisimple factors.  Indeed,
it permutes the simple factors of $E/Z(E)$, and each $E_i$ is
the derived subgroup of the preimage of its simple factor, so
the perfect lifts are intrinsic.  The matrix algebras generated
by the $E_i$ are the full factor algebras
$\operatorname{End}(V_i)\otimes1$.  A normalizing matrix therefore
acts as a product of linear factor maps followed by the induced
permutation.  This follows by the matrix-unit proof that every
automorphism of a full matrix algebra is inner, followed by its
tensor-factor centralizer; no rational field-of-definition
assertion is needed.

First suppose $g\in J$ fixes every tensor factor.  Then
$g=A_1\otimes\cdots\otimes A_t$, with each $A_i$ normalizing
the represented $E_i$.  If $g$ is nonscalar, some $A_i$ is
nonscalar.  D-SPOR-02 bounds its largest eigenspace by $3a_i/4$.
For tensor products of diagonalizable operators, the largest
eigenspace multiplicity is at most the largest multiplicity in
one factor times the product of the other dimensions: sum over
the eigenspaces in the other factors.  Thus
\[
                              \nu(g)\ge\dim V/4.             \tag{SPOR.5}
\]
The factor operators are diagonalizable up to scalar since $g$
has finite order; alternatively D-SPOR-02 represents each one
as a scalar multiple of a finite-order matrix.

Now suppose $g$ has a tensor-factor cycle of length $l\ge2$.
All factors along that cycle have a common dimension $a$.
On their tensor product the operator has the form of cyclic
rotation composed with invertible maps between successive factors.
Transporting bases around the cycle conjugates it by a product
of changes of basis to cyclic rotation with a single remaining
twist $B$ on one factor.  This twist is diagonalizable.  To see
this without a unitarity convention, a power of the full $g$ is
the identity, so a power of each cycle operator is scalar.
Its $l$-th power is conjugate to $B^{\otimes l}$, and a power
of $B$ must therefore be scalar.  The polynomial $x^k-c$ has
distinct roots in characteristic zero when $c\ne0$.

Choose a diagonalizing basis for $B$.  The basis words of length
$l$ on an alphabet of size $a$ are then permuted by cyclic
rotation, with nonzero scalar weights.  On an orbit of length
$d$, the operator has characteristic polynomial $x^d-c$ with
$c\ne0$, so any specified eigenvalue occurs at most once in
that orbit.  There are exactly $a$ constant words; every other
word orbit has length at least two.  Therefore the largest
eigenspace on this cycle has dimension at most
\[
       a+\frac{a^l-a}{2}=\frac{a^l+a}{2}
                  \le\frac34 a^l,                         \tag{SPOR.6}
\]
because $a^{l-1}\ge2$.  Eigenvalue coincidences between different
orbits do not change this count.  Tensoring with the other cycle
operators gives (SPOR.5) for the full $g$ as well.

We now apply the bound to the actual intermediate cover attached
to the inducing representation.  Put $e=\dim V$ and let $h$ be
the genus of its base.  Its multiplicity cohomology has complex
dimension six, so with genuine inertias $g_1,\ldots,g_r$,
\[
       6=2e(h-1)+\sum_j\operatorname{rank}(1-g_j).           \tag{SPOR.7}
\]
Frobenius reciprocity identifies this complex Hodge filtration
with the original one.  Its conjugate is determined by the actual
Hermitian polarization, so the period rank factors through this
intermediate cover even for a larger inducing character field.
Finite covers of a fixed pointed curve have only finitely many
topological possibilities.  If $h=1$, (SPOR.7) gives at most six
marked points and period rank at most six; if $h\ge2$, it is
impossible because $e\ge6$.  Rank nine therefore forces
\[
               h=0,\qquad r\ge12,\qquad
               \sum_j\operatorname{rank}(1-g_j)=2e+6.
                                                               \tag{SPOR.8}
\]
For every nonscalar inertia, (SPOR.5) gives actual moving rank
at least $\lceil e/4\rceil$.  A genuine scalar inertia has moving
rank $e$, so obeys the same lower bound.  Hence
\[
                        12\lceil e/4\rceil\le2e+6.         \tag{SPOR.9}
\]
This is impossible for $e>6$, since its left side is at least
$3e$; for $e=6$ its left side is $24$ and its right side $18$.
Finally $e\ge6$ follows from (SPOR.4).  The contradiction proves
the theorem.  In particular it is not merely an obstruction to
one branch profile or one small sporadic representation.
\end{proof}

\paragraph{Reproduction and scope.}
The complete group tables are an explicit finite published input.
The two independent arithmetic implementations verify the spectral
inequality without floating-point computations.  They do not
reprove the identification of ATLAS tables with their groups or
the classification of sporadic groups.  The theorem assumes the
normal central product acts irreducibly; a reducible normal subgroup
is not silently substituted.  D-BI-19 supplies that precise hypothesis
for the primitive tensor-indecomposable blocks to which it applies.

